\documentclass{beamer}
\usepackage{animate}
\usepackage{amsmath}
\usepackage{tikz}
\usetheme{Madrid}
\usepackage{tikz}
\usepackage{soul}
\usepackage[normalem]{ulem}
\graphicspath{{./figures/}}

\useinnertheme{circles}
%Information to be included in the title page:
\title{Maximum likelihood estimation (MLE) recap}
\author{Jakub Koubele}
\date{13th April, 2026}

\renewcommand{\insertshortauthor}{Jakub Koubele}  
\renewcommand{\insertshorttitle}{Maximum likelihood estimation}
\renewcommand{\insertshortdate}{13th April, 2026} % Define custom text for 3rd block

\usepackage{amsthm}
\newcommand*{\QED}{\hfill\ensuremath{\blacksquare}}





\begin{document}
	\frame{\titlepage}
	
	
	\begin{frame}
		\frametitle{Introduction}	
		\pause		
		\begin{itemize}
		\item If we have some \textbf{data}, we may want to fit a \textbf{model} on them. \pause
		\item Broadly, our motivation for modeling the data may be:
			\begin{itemize}
				\item We may want to \textbf{predict} future/hidden outcomes (e.g., diagnose a patient based on symptoms and biomarkers).
				\item We may want to \textbf{infer} causality/associations (e.g., how lifestyle choices influence lifespan).
			\end{itemize}
			
		\pause
		\item Many models are \textbf{probability distributions}.
		\begin{itemize}
			\item These may be 'simple' distributions (e.g., normal distribution), different sort of (generalized-) linear models, artificial neural networks, etc.
		\end{itemize}
		\pause
		\item Probability distribution models are most commonly fit by \textbf{maximum likelihood estimation} (MLE).
		
		
		\end{itemize}		
	\end{frame}	
	
	
	\begin{frame}
		\frametitle{Gaussian (normal) distribution}
		\begin{itemize}
			\item Consider the Gaussian distribution, given by a density function
			$$
			f(x) = \frac{1}{\sqrt{2 \pi \sigma^2}} \exp \left( - \frac{\left( x-\mu \right)^2 }{2\sigma^2} \right)
			$$
			where $\mu \in \mathbb{R}, \sigma^2 \in \mathbb{R}_+$ are \textbf{parameters}.
			\pause
			\item We are given some data e.g., measurements of physical height. E.g. we may have $n=5$ observations like this:
			$$
			X = \{ 161\text{ cm}, 193\text{ cm}, 147\text{ cm}, 181\text{ cm}, 175\text{ cm}\}
			$$
			\pause
			\item Our task is to estimate the parameters from the data.
			
		\end{itemize}
	\end{frame}	

\begin{frame}
	\frametitle{Likelihood function}
	\begin{itemize}
		\item Recall the definition of likelihood function:
		$$
		L(\mu, \sigma^2) = \prod_{i=1}^{n} f(x_i \mid \mu, \sigma^2)
		$$
		\pause
		\item We may plot the likelihood (for a specific data $X$ from the example above):	
		
			
	\end{itemize}
\vspace{0.3cm}

\begin{columns}
	\column{0.5\textwidth}
	\centering
	\includegraphics[width=0.95\linewidth]{likelihood_full.png}
	
	\column{0.5\textwidth}
	\centering
	\includegraphics[width=0.95\linewidth]{likelihood.png}
\end{columns}
\end{frame}	

\begin{frame}
	\frametitle{Log-likelihood}
	\begin{itemize}
		\item Optimizing the likelihood function directly is cumbersome, due to the product term. Fortunately, for all positive functions $f(x)$, it is
		$$
		\arg \max f(x) = \arg \max \log(f (x))
		$$
		\pause
		\item So instead of maximizing likelihood, we may equivalently maximize log-likelihood, given as:
				
		\begin{align*}
			l(\mu, \sigma^2)
			&=\log \left( L(\mu, \sigma^2) \right)
			=\log \left( \prod_{i=1}^{n} f(x_i \mid \mu, \sigma^2) \right)			 \\
			&= \sum_{i=1}^{n} \log f(x_i \mid \mu, \sigma^2)
		\end{align*}
		
\end{itemize}		
	
\end{frame}	

\begin{frame}
	\frametitle{Log-likelihood: visualization}
	
	\centering
	\includegraphics[width=0.75\linewidth]{log_likelihood.png}	
	
\end{frame}	

\begin{frame}
	\frametitle{Solving for $\hat{\mu}$}
	
	
	
	
	\pause
	\begin{itemize}
		\item We have	
		\begin{align*}
			l(\mu, \sigma^2)
			&= \log \left( \prod_{i=1}^{n} f(x_i \mid \mu, \sigma^2) \right)	\\
			&= -\frac{n}{2}\log(2\pi\sigma^2)
			- \frac{1}{2\sigma^2} \sum_{i=1}^{n} (x_i - \mu)^2
		\end{align*}
	\pause
	\item Differentiating with respect to $\mu$:
	$$
	\frac{\partial l(\mu, \sigma^2)}{\partial \mu} = \frac{1}{\sigma^2} \sum_{i=1}^{n} (x_i - \mu)
	$$	
	\pause
	\item By setting the derivative to 0 and solving for $\mu$, we get
	$$
	\hat{\mu} = \frac{1}{n} \sum_{i=1}^{n} x_i
	$$		
		
	\end{itemize}
	
		
	
\end{frame}	

	
	
	
\end{document}
